% =============================================================================
% HelloWorld.tex — Minimal Beamer sample for verifying the workflow
%
% This deck uses Beamer's built-in Madrid theme and has NO dependency on
% Preambles/header.tex. It compiles on a fresh fork without any setup.
%
% DELETE THIS FILE and replace with your own lectures once you're ready.
%
% Compile with: /compile-latex HelloWorld
% Or manually:  cd Slides && xelatex HelloWorld.tex
% =============================================================================

\documentclass{beamer}
\usetheme{Madrid}

\title{HelloWorld: Academic Workflow Sample}
\author{Your Name}
\institute{Your Institution}
\date{\today}

\begin{document}

\frame{\titlepage}

\begin{frame}{Why This Sample Exists}
    \begin{itemize}
        \item Verifies that XeLaTeX is installed and working
        \item Confirms \texttt{/compile-latex} produces a valid PDF
        \item Exercises the full 3-pass + bibtex pipeline (see citation below)
        \item Provides a starting point before you build real lectures
    \end{itemize}
    Placeholder citation: see \cite{Example2024_book}.
\end{frame}

\begin{frame}{Next Steps}
    \begin{enumerate}
        \item Compile this deck: \texttt{/compile-latex HelloWorld}
        \item Translate to Quarto: \texttt{/translate-to-quarto HelloWorld.tex}
        \item Delete this file and create your own lecture
    \end{enumerate}
\end{frame}

\begin{frame}{References}
    \bibliographystyle{plain}
    \bibliography{../Bibliography_base}
\end{frame}

\end{document}
